\section{Background}
\label{sec:background}

% TODO: Write background


\subsection{Hardware-Software Interfaces in Tensor Accelerators}
\label{subsec:hw-sw-interface}

Tensor accelerators occupy a fundamentally different point in the design space than scalar or vector processors.
Rather than executing fine-grained instruction streams over flat register files and cache-coherent memory, they expose \emph{coarse-grained} operations that consume and produce multi-dimensional tensors and replace implicit caches with \emph{explicitly managed} on-chip memories (scratchpads, accumulator banks).
Gemmini~\cite{gemmini-dac2021}, Berkeley's open-source systolic-array generator, is representative: it attaches to a Rocket core as a RoCC accelerator~\cite{chipyard} with non-standard RISC-V custom instructions, and its programmer-visible state consists of a row-addressed scratchpad and an accumulator.
Its ISA exposes only configuration, DMA data-movement, and matrix-multiply execution instructions---every operation simultaneously specifies \emph{what} to compute, \emph{where} operands live in a private memory hierarchy, and \emph{how} data should be staged.

This distinctiveness places heavy demands on compiler infrastructure.
ML frameworks feed tensor IR (StableHLO, Linalg, TOSA) into backends such as XLA~\cite{xla}, TVM~\cite{tvm}, or IREE~\cite{iree}, which must lower high-level operations to the accelerator's native instructions.
Doing so correctly requires tiling and fusing loop nests to fit scratchpad capacity, scheduling explicit DMA transfers, allocating scratchpad regions, and orchestrating double-buffering---all of which presuppose a precise model of what each instruction does to which piece of tensor state.

Consider lowering a single XLA \texttt{dot} $C{=}A{\cdot}B$ for matrices exceeding Gemmini's $16{\times}16$ tile.
The backend must emit a tiled sequence: \texttt{config\_ex}/\texttt{config\_ld}/\texttt{config\_st} to set dataflow and stride parameters; \texttt{mvin}/\texttt{mvin2} to stage tiles of $A$ and $B$ into the scratchpad; \texttt{preload} followed by \texttt{compute\_preloaded} (or \texttt{compute\_accumulated}) to fire the systolic array for each output tile; and \texttt{mvout} to drain results.
Generating this sequence requires the backend to know, for every instruction, which tensor tile it touches, which scratchpad region it occupies, and how its output composes with subsequent operations.

This is precisely where most accelerators fall short.
Mapping tensor IR to an instruction sequence presupposes a \emph{machine-readable} ISA that specifies both the binary encoding and the \emph{tensor-level semantics} of each instruction.
For Gemmini, semantics live informally in prose documentation, the Chisel RTL, and a hand-written Spike functional model; most open-source accelerators offer only RTL and perhaps a driver.
Building a backend today is therefore largely an exercise in \emph{manual reverse-engineering of the RTL}---a human infers what each custom instruction means at the tensor level and re-encodes that knowledge by hand.
This is the bottleneck that motivates automatically lifting tensor-level ISA semantics directly from RTL.

\subsection{Automated Compiler Generation and the Specification Bottleneck}


\subsection{RTL Extraction and MLIR-Based Semantic Lifting}
